\begin{abstract}
7777Put your abstract here. Abstracts typically start with a sentence motivating why the subject is interesting. Then mention the data, methodology or methods you are working with, and describe results. 
\end{abstract}

\section{Introduction}\label{sec:intro}

deutsche Medaillenkrise bei den Olympischen Sommerspielen 2024 

7777Motivate the problem, situation or topic you decided to work on. Describe why it matters (is it of societal, economic, scientific value?). Outline the rest of the paper (use references, e.g.~to \Cref{sec:methods}: What kind of data you are working with, how you analyse it, and what kind of conclusion you reached. The point of the introduction is to make the reader want to read the rest of the paper.

\section{Data and Methods}\label{sec:methods}

For the analysis of performance evolution and performance density evolution we made some restrictions to the dataset to have comparable and sufficient data.
We Analysed for the U18, U20 and Adult age groups, representing the most relevant age groups witth consistent data. U23 didn't had enough data, also for younger age groups.
Also always used the top 35 Ranks for analysis to have enough data for all disciplines and age groups.
A Group is defined as a combinaion of Gender x Age group x Discipline to ensure consistency with the original leaderboards.
We compared two time periods: Data Start period (2001-2004) and End data period(2021-2024)
We used four year periods to exlude one-year variances.

To measure the overall performance improvement of a group, we use the Median Difference ($\Delta \tilde{x}_{\%}$).
The median of each period ($p$) is calculated from the respective performance pools ($S_p$) of the four year periods.
The relative change between the start ($\tilde{x}_s$) and the end period ($\tilde{x}_e$) is defined as:

\begin{equation}
\Delta \tilde{x}_{\%} = \left( \frac{\tilde{x}_e - \tilde{x}_s}{\tilde{x}_s} \right) \times 100
\end{equation}

A positive $\Delta \tilde{x}_{\%}$ indicates an overall improvement in performance, while a negative $\Delta \tilde{x}_{\%}$ shows a decrease in performance scores.
Statistical significance for the comparison of the medians is measured using the Mann-Witney U test.
Without assuming a normal distribution, this test allows to evaluate whether the distributions of the two periods differ significantly.

To measure \textbf{Performance Density (PD)}, we used the performance spread ($G_p$) from annual top 3 performances to the bottom 3 performances (e.g. Rank 33-35) of our dataset for the two time periods.
Mathematically, the gap ($G_p$) for a period ($p$) is defined as the difference between the collective mean of the top 3 and bottom 3 performances recorded across all years ($n$) within that period.

\begin{equation}
G_p = \frac{1}{n}\sum_{t=1}^{n} \bar{x}_{t, \text{top3}} - \frac{1}{n}\sum_{t=1}^{n} \bar{x}_{t, \text{low3}}
\end{equation}

Furthermore, the evolution of the performance gap ($\Delta G_{\%}$) between the start period ($G_s$) and end period ($G_e$) is calculated as

\begin{equation}
\Delta G_{\%} = \left( \frac{G_e - G_s}{G_s} \right) \times 100
\end{equation}

A positive $\Delta G_{\%}$ suggests a decrease in \textbf{PD}, while a negative $\Delta G_{\%}$ suggests a increase in \textbf{PD}.

To assess the statistical significance of the change in \textbf{PD}, we use a bootstrap resampling method with $B=10,000$. 
We draw with replacement from the performance pools, we generated a distribution of the difference in gaps ($\Delta G_{obs}$). 
We then calculate the $95\%$ confidence interval (CI) for the difference in gaps. 
If the CI includes zero, the null hypothesis of no significance change in \textbf{PD} cannot be rejected.


FOR RANK Analysis:

- Calculation of the mean per rank per group per period
- calculation of the difference of the two periods per rank per group

% This is the template for a figure from the original ICML submission pack. In lecture 10 we will discuss plotting in detail.
% Refer to this lecture on how to include figures in this text.
% 
% \begin{figure}[ht]
% \vskip 0.2in
% \begin{center}
% \centerline{\includegraphics[width=\columnwidth]{icml_numpapers}}
% \caption{Historical locations and number of accepted papers for International
% Machine Learning Conferences (ICML 1993 -- ICML 2008) and International
% Workshops on Machine Learning (ML 1988 -- ML 1992). At the time this figure was
% produced, the number of accepted papers for ICML 2008 was unknown and instead
% estimated.}
% \label{icml-historical}
% \end{center}
% \vskip -0.2in
% \end{figure}

\section{Results}\label{sec:results}

PLOT: MEDIAN vs GAP Difference START vs END Period (with gender color and age group markers), Significance for different 

Generally a trend that there are more groups with a significant higher PD (n=) to n=x lower PD. 
For the median this is more a little more balanced (sig. higher med n= to lower med n=)

For the general median trend, we see that in the women's adults 100\% of their median increases are statistically significant.
For the men its 81.8\%.
In U18 and U20 Men's there are more decreases in median performances, for womens its more or less half improvment and decreasement.

The clearest trend is seen in the U20 Men's category, where X\% of the groups showed a significant Polarization of performance.
This is followed by the U18 Men's category, where X of the groups showed a significant Polarization of performance.
Adult Men and Women both displayed significant decrease in PD on over X\% of the groups (X and X respectively)
In all discipline groups around the same amount of decrease in PD is observed.
A Compression was not often seen and mostly non-significant across all groups.

PLOT: RANK Evolution

    - PLOT shows the rising discrepance: first places are getting better, lower places worse
    - In general, but especially the M U18 and U20, while women -> better over all places

PLOT: YEAR EVOLUTION/ PERIOD EVOLUTION

    - Men: Decrease in Mean over time -> sign. drop of 1.4 \% from 2013-16 to 2017-2020
    - Std decrease step by step around 5\% per period, 2005-08 dirst improvment, 
    - Women: step by step increase in mean
    - STD deviation just in the last period 2021-24 sign. higher


\section{Discussion \& Conclusion}\label{sec:conclusion}

- As the results show, in general there is a problem with a rising discrepance of top performers to "second row"
- More seen in Men than in women, and more in younger age groups, in line with the "Nachwuchs problem?"
- More problems in technical disciplines -> Corona, less people in the sport -> less trainer, techn. needs more attention
- Running marathon strong improvments -> shoes etc
- Men Step by step decrease MEan and std -> also a reason of support?
- Women getting better: later start in the sport, gender discrepancy, improvment, evolution still possible

\newpage

\section*{Contribution Statement}
Explain here, in one sentence per person, what each group member contributed. For example, you could write: Max Mustermann collected and prepared data. Gabi Musterfrau and John Doe performed the data analysis. Jane Doe produced visualizations. All authors will jointly wrote the text of the report. Note that you, as a group, a collectively responsible for the report. Your contributions should be roughly equal in amount and difficulty.

\section*{Notes} 

Your entire report has a \textbf{hard page limit of 4 pages} excluding references and the contribution statement. (I.e. any pages beyond page 4 must only contain the contribution statement and references). Appendices are \emph{not} possible. But you can put additional material, like interactive visualizations or videos, on a githunb repo (use \href{https://github.com/pnkraemer/tueplots}{links} in your pdf to refer to them). Each report has to contain \textbf{at least three plots or visualizations}, and \textbf{cite at least two references}. More details about how to prepare the report, inclucing how to produce plots, cite correctly, and how to ideally structure your github repo, will be discussed in the lecture, where a rubric for the evaluation will also be provided.